Studiamo lo spettro di emissione $ \gamma $ di varie sorgenti radioattive come \ce{^60Co}, \ce{^137Cs}, \ce{^22Na} etc. misurando la radiazione ocn uno scintillatore del tipo \ce{NaI(Tl)}. Osservando la posizione dei picchi decretiamo il processo che origina l`emissione osservata, caratterizziamo la presenza di elementi radioattivi in una roccia di origine ignota e lo spettro emissivo del fondo ambientale. Studiamo la risoluzione del rivelatore al variare dell`energia osservata stimandola a $ R = 6.5\% $, e successivamente misuriamo l`attività di una sorgente di tipo \ce{^137Ce} con metodo relativo e assoluto. Infine determiniamo il coefficiente di assorbimento di massa $ \mu / \rho = (0.1040 ± 0.0010)$ \unit{cm^2/g} compatibile al $ 5\% $ con il valore teorico considerato.
